% 本文件由 LaTeX 排版 Agent 根据有效 translation.json 自动生成。
% author=Miscellaneous; work_id=0863; title=玛拉，塞拉皮雍之子，书信

\chapter{\WorkTitle{\Person{玛拉}，\Person{塞拉皮雍}之子，书信}}

% structure: p0001 source=h1 target=omitted reason=page-title-already-rendered-by-parent

\Person{玛拉}，\Person{塞拉皮雍}之子，致我儿\Person{塞拉皮雍}：愿你平安。\par

当你的师长与监护人写信给我，告知我你\Emph{虽}年纪尚幼，却勤于学业时，我感谢天主，因你一个小小童子，在无人\Emph{指引}之下，竟能如此认真开始；而对我自己而言，这\Emph{也}是一种安慰——因我听闻你，\Emph{虽为幼童}，却\Emph{已显露出}如此伟大的心胸与良知；\Emph{这种品格}，在许多人身上，虽\Emph{曾良好开端}，却未见有继续的热忱。\par

为此，看，我写下了这篇记录，\Emph{涉及}我通过细心观察在世上所发现的一切。因为人生的样式已为我所细察。我行走于学问之路，从希腊哲学的研究中发现了这一切，尽管它们在生命诞生之时遭遇了船难。\par

所以，我儿，当勤勉\Emph{专注于}那些合乎自由人之事，\Emph{以便}投身学问，追随智慧；并努力在已开始的\Emph{习惯}中坚定自己。也要记念我的训诫，如同一个安静而热爱学问追求的人。即使\Emph{这样的生活}在你看来十分艰苦，\Emph{然而}当你稍加体验之后，它便会变得非常甘美：因为我自己也曾如此。再者，当一个人离开家乡，仍能\Emph{依旧}保持\Emph{从前的}品性，并妥善尽其所当为，他便是那蒙拣选之人，被称为\Quote{天主的祝福}，且找不到任何可与他的自由相比之物。因为那些被召于追求学问的人，是在寻求摆脱时世的纷扰；那些把握智慧的人，是紧握义德的盼望；那些立足于真理的人，是展现其德行的旗帜；那些修习哲学的人，是盼望逃离世间的烦恼。因此，我儿，你也要在这些事上明智行事，如同一个寻求度纯净生活的智慧之人；\Emph{当心}，免得众人所贪图的利益使你软弱，你的心思转向贪求没有稳固的财富。因为，财富若以欺诈得来，必不持久；即便公正\Emph{取得}，也不能长存；而你在世上所见的一切，既属于\Emph{仅}短暂之物，\Emph{都注定}要如梦幻般逝去：它们\Emph{不过如同}季节的起落。\par

至于人生所充满的\Emph{那种}虚荣\Emph{对象}，也不必忧虑：因为那些给予我们快乐之事，往往迅速带来伤害。\Emph{最明显不过的是}心爱子女的诞生。因为它在两方面显然带给我们伤害：对于有德者，\Emph{我们}对他们的情感折磨我们，因他们\Emph{卓越的}品性我们受煎熬；对于堕落者，我们为管教他们而烦恼，因他们的恶行而痛苦。\par

此外，你已听闻关于我们同伴的事：当他们离开\Place{萨莫萨塔}时，\Emph{对此}感到悲伤，\Emph{仿佛}埋怨\Emph{所处的}时代，如此说道：\Quote{我们现在远离家乡，不能再返回我们的城市，或看见我们的族人，或向我们的诸神献上赞美的问候。}那日理应被称为\Emph{哀哭}之日，因为一种沉重的悲伤同样占据了他们众人。因为他们记念自己的父亲而哭泣，\Emph{想起}母亲而呜咽，为兄弟而忧伤，\Emph{并为}留在身后的未婚妻而悲痛。尽管我们听说他们从前的同伴正前往\Place{塞琉西亚}，我们还是暗中\Emph{启程，}上路朝他们而去，将自己的苦难与他们的相连。于是我们的悲伤极其猛烈，我们的哭泣理当涌流，因我们绝望的处境，我们的哀号\Emph{聚集成}浓密的乌云，我们的苦难比山更重：因为我们中没有一人能够抵挡那袭击他的灾祸。因为对活人的爱是强烈的，对死者的悲伤也是深重的，而我们的苦难正驱使我们无路\Emph{可逃}。因为我们看见我们的兄弟和儿女成了俘虏，我们记念已故的同伴，他们被安葬在异乡。我们每人也都为自己焦虑，唯恐祸不单行，或新的灾难接续前一个。身为囚徒、\Emph{且}经历\Emph{这般}遭遇的人，能有什么享受呢？\par

至于你，我所爱的，不要因你孤身一人而被逐往各处而忧伤。因为人本是为此而生，\Emph{注定}要遭遇时世的变迁。\Emph{反而}你当这样想：对智慧人而言，处处皆是故乡；\Emph{而且}在每座城中，善人有许多父亲和母亲。\Emph{你若不信}，便从\Emph{你亲眼所见}的事上取一个证据。有多少不认识你的人爱你如同\Emph{自己的}儿女；\Emph{又有}多少妇人接待你如同\Emph{对待}自己的亲人！实在，你身为异乡人却蒙了眷顾；实在，因你小小的爱心，许多人已对你生出炽热的情感。\par

关于盘踞在世上的谬误，我们还能说什么呢？由于劳苦，人生之旅令人痛苦，因它的动荡，我们像风中的芦苇，时而弯向这边，时而弯向那边。因为我曾惊讶于许多抛弃自己儿女的人，也惊诧于另一些人抚养并非己出的孩子。有人获取世上的财富，我也曾惊诧于另一些人继承并非自己\Emph{所得}的东西。这样\Emph{你可}以明白并看到，我们是在谬误的引导下行走。\newline{}\par

最智慧的人啊，请开始告诉我们，一个人究竟能依靠他的哪样\Emph{财产}，或者说关于哪些事他能说它们是恒久的？\Emph{你会说是}丰富的财富吗？它们会被抢走。堡垒吗？它们会被摧毁。城市吗？它们会变成废墟。尊荣吗？它会被贬抑。华贵吗？它会被倾覆。美貌吗？它会凋谢。或是法律吗？它们会失效。或是贫穷吗？它受人轻视。或是儿女吗？他们会死去。或是朋友吗？他们会背信。或\Emph{是人的}赞誉吗？嫉妒先于它们。\newline{}\par

因此，让人因自己的帝国而欢喜，如达里乌斯；或因好运，如波利克拉特斯；或因勇武，如阿喀琉斯；或因妻子，如阿伽门农；或因后代，如普里阿摩斯；或因技艺，如阿基米德；或因智慧，如苏格拉底；或因学问，如毕达哥拉斯；或因才智，如帕拉墨得斯——我儿，人的生命离开世界，但他们的赞誉与德行永存。\newline{}\par

那么，我的小儿，选择那不褪灭的。因为专心于此的人被称为谦逊，\Emph{且}为人所爱，是美名的爱慕者。\newline{}\par

再者，若有不幸之事临到你，不要归咎于人，也不要对天主发怒，也不可咒骂你所处的时代。\newline{}\par

你若持守此心，你从天主所得的赏赐就不小，它无需财富，也永不陷入贫穷。因为你要毫无惧怕地度过一生，且满心喜乐。因为恐惧为\Emph{自己}本性辩护，不属于智者，而属于悖逆律法的人。因为从没有人被夺去智慧，如同被夺去财产一样。\newline{}\par

要殷勤追随学问，胜过财富。因为\Emph{一个人的}财产越多，\Emph{伴随它们的}祸患也越大。因为我曾亲眼看见，\Emph{一个人的}财物多，临\Emph{到他}的苦难也多；奢侈聚集之处，忧伤也汇集；财富丰盈之处，\Emph{积蓄着}多年的苦涩。\newline{}\par

因此，你若有智慧行事，并殷勤关注\Emph{你的行为}，天主必不停止帮助你，人也必不停止爱你。\newline{}\par

让你所能获得的东西满足你；此外，你若能一无所有，你就有福了，也没有任何人会嫉妒你。\newline{}\par

也请记住这一点：除了\Emph{贪爱}利益，没有什么会极大地扰乱你的生命；而且，死后无人被称为产业的主人。因为正是由于这种欲望，软弱的人被掳掠，他们不知道一个人住在自己的财物中，\Emph{只是}像一个偶然的过客，他们终日恐惧，因为\Emph{这些财物}并不属于他们：因为他们抛弃了属于自己的，而去追求不属于他们的。\newline{}\par

当智者被暴君的手强行拖走，他们的智慧因诽谤而被剥夺自由，他们因\Emph{卓越的}才智而遭掠夺，却\Emph{没有辩护的机会}时，我们还能说什么呢？\Emph{他们并非完全值得怜悯}。因为雅典人处死苏格拉底得到了什么益处？他们得到的报应是饥荒和瘟疫。萨摩斯人焚烧毕达哥拉斯得到了什么益处？在一小时内，他们的整个国家被沙土覆盖。犹太人\Emph{杀害}他们的智王又得到了什么益处？从那时起，他们的王国就被\Emph{从他们手中}夺走了。因为天主公义地赐予\Emph{所有}三位智者的智慧以报应。雅典人死于饥荒；萨摩斯人被海水不可挽回地淹没；犹太人则荒凉，被逐出王国，流散到各地。\Emph{不}，苏格拉底\Quote{没有}死，因为柏拉图；毕达哥拉斯也没有，因为赫拉雕像；智王也没有，因为他所颁布的新法律。\newline{}\par

此外，我儿，我仔细观察了人类，他们陷于何等悲惨的境地。我惊讶他们没有被周围的灾祸完全击倒，\Emph{而且}甚至\Emph{他们的}战争、\Emph{所受的}痛苦、疾病、死亡、贫穷都不够他们受的；\Emph{而且}他们像野兽一样，在\Emph{彼此的}敌意中必须互相冲撞，\Emph{试图}看谁能给同伴造成更大的伤害。因为他们脱离了真理的界限，违背了一切诚实的法律，只求满足私欲；因为当一个人急切\Emph{获取}他所欲求的东西时，他怎么可能适当地做他应当\Emph{做的事}呢？他们不承认任何约束，很少向真理和良善伸出手，在生活中却像聋子和瞎子一样行事。而且，恶人欢喜，义人不安。有\Emph{的}说没有，没有的竭力去获取。穷人寻求\Emph{帮助}，富人隐藏\Emph{财富}，人人彼此嘲笑。醉酒的人麻木，清醒的人羞愧。有人哭泣，有人歌唱；有人欢笑，有人忧愁。他们以恶事为乐，厌恶说实话的人。\newline{}\par

那么，当世人欲以世人\Emph{自己的}蔑视来摧折人，既然他们\Emph{和祂}没有\Emph{同一}生活方式，人又何必惊讶呢？\Quote{这些}是他们所关切的。其中一人正在\Emph{盼望那时刻}，当他在战中获得胜利的荣名；然而勇敢的人竟不知有多少愚蠢的欲念将人掳获于世间。但愿悔改之心能短暂地临于他们！因为，他们虽因勇武而获胜，却被贪欲之力所征服。我曾试验过人，结果如此：他们所专注的唯有财富的丰裕。因此，他们也没有确定的目标；因着心意的不定，人突然\Emph{从精神高昂}跌入忧伤的吞噬。他们不注视永恒的浩瀚财富，\Emph{也不思想}每一次患难的造访都同样将我们引向同一\Emph{最终}终点。因为他们委身于肚腹的尊严——\Emph{那}恶人德行上\Emph{的}巨大污点。\par

此外，\Emph{关于}这封\Emph{信}我意欲写给您的信，阅读它还不够，最好的乃是付诸实行。因我自知，当您亲身体验这种生活方式后，它必将令您十分愉悦，且您将免于剧痛；因为我们忍受财富，\Emph{仅仅}是为了儿女。\par

因此，\Emph{最}受爱戴的人啊，请将忧愁从您身上除去——那对\Emph{人}绝无益处之物；并将挂虑驱走——那丝毫不带来益处之物。因为我们没有\Emph{可资利用的}资源或计谋——\Emph{唯}有一颗\Emph{能够}应付灾难、承受我们时常从时代领受的磨难的心。因为我们理当注视这些，而非\Emph{仅}那些充满喜乐与美名之事。\par

当您致力于智慧——万善之泉、不朽之宝。您将安息其上，得享安宁。因为智慧将成为您的父与母，以及您人生中的良伴。\par

请与刚毅和忍耐——那些\Emph{德行}，能\Emph{成功}应对降临于软弱之人身上的患难——结成至亲之交。因其力量如此强大，足以承受饥饿，\Emph{能}忍受干渴，并缓解一切痛苦。对于劳苦，甚至对于解体，它们亦能欣然以对。\par

请殷勤留意这些事，您将度无忧无虑的生活，我也将得安慰，而您将被称爲\Quote{父母的喜悦}。\par

因为在往昔，当我们的城邑仍在其伟大之中屹立时，您或许知道，有许多人遭致可憎的言语攻击\Emph{在我们中间}；但就我们自己而言，我们早已承认，我们从她的众多子民那里接受了爱与尊荣，不亚于任何人：只是时代的状态\Emph{唯独}阻止了\Emph{我们的}完成所决定的事。而在此监牢中，我们感谢天主，因为我们得到了许多人的爱：因我们竭力维持节制与愉快的生活；若有人强行驱赶我们，他\Emph{不过}是公开作证反对自己，证明他疏离一切善事，且他将因自身的污点——\Emph{在他身上的}——而蒙受羞耻与羞辱。因为我们已显明了我们的真理——\Emph{那真理}，在我们\Emph{如今}倾覆的王国中不曾拥有的真理。但若罗马人本着公义与正直，允许我们\Emph{回去}到本国——\Emph{被吁请}\Emph{去做}——他们便显为仁慈之人，将赢得良善与公义之名，同时亦\Emph{将拥有}一个安居的和平国度：因为他们若使我们得自由，便将彰显其伟大；\Emph{并且}我们将顺从时代所赋予的主权。但不要让他们像暴君一样，驱赶我们\Emph{仿佛我们是}奴隶。然而，若事情已\Emph{已经}注定，我们将领受的无非是那为我们存留的平安之死——\Emph{可怕的}死亡。\par

但你，我的小儿，若你决意殷勤了解这些事，首先须节制食欲，对你所\Emph{纵情}之处设定界限。寻求抑制怒气的能力；不要\Emph{顺从}激情的爆发，而要聆听良善的\Emph{催促}。\par

至于我，我从今以后所挂虑的是\Emph{这一点}：尽我所忆\Emph{从前的}，留下一部\Emph{记载它们}的书，并以审慎之心完成我被指派\Emph{要走}的路程，无苦痛地离开这世间的悲伤苦难。因我的祈祷是：愿我得\Emph{我的}释；至于\Emph{通过}何种死亡，我毫不介意。但若有人\Emph{为此}感到困扰或焦虑，我无法给他建议：因为在那边，在普世的居所，他将发现我们已在他之前。\par

他的一位朋友，在锁链中坐在\Person{玛拉·塞拉比翁}旁边，问他：\Quote{不，我以你的生命起誓，玛拉，告诉我你看见了什么\Emph{缘由}使你发笑？}\Quote{我在笑，}玛拉说，\Quote{因时间虽未从我这里借取任何邪恶，却正在偿还我。}\par

\Emph{这里}结束\Person{玛拉·塞拉比翁}的书信。\par

\SourceRecord{Translated by B.P. Pratten.；Ante-Nicene Fathers；Vol. 8.；Edited by Alexander Roberts, James Donaldson, and A. Cleveland Coxe.；Buffalo, NY: Christian Literature Publishing Co.,；1886.；<http://www.newadvent.org/fathers/0863.htm>.}
